
\section{Weak coupling: what a proof of the mass gap still needs}\label{sec:weak}

The continuum limit of four-dimensional Yang--Mills theory is taken along the
tuning dictated by asymptotic freedom: for $\SU(N)$ the bare coupling
$g^2=2N/\beta_W$ tends to zero like $1/(b_0\log(1/(a\Lambda)))$ with
$b_0=11N/(48\pi^2)$ and a fixed physical scale $\Lambda$, so $\beta_W\to\infty$
logarithmically in $1/a$. Theorems~\ref{thm:A} and~\ref{thm:bpstrong} hold for
$\beta_W<N^2/29$ and $\beta_W<N^2/42$ respectively; at fixed $\beta$ the lattice
gap in physical units, $m_G(\beta)/a$, diverges as $a\to0$, and no continuum
limit exists at fixed $\beta$. The results of this paper therefore say nothing
about the continuum theory without statement (D). Three facts show what (D)
asks for.

\emph{(a) A cube condition is a mass.} By Theorem~\ref{thm:free} the best
constant for a free field is at most $s^2/(s^2+16d/9)$ with $s=m\ell$. At weak
coupling the leading approximation of the Wilson measure, in any gauge, is a
massless Gaussian field, for which the constant vanishes as $a\to0$ at fixed
$\ell$. Hence (D) does not follow from perturbation theory to any finite order:
it holds only if the interaction generates a mass of order $1/\ell$, and it is
a quantitative form of that mass generation.

\emph{(b) The cube must grow.} For $\SU(N)$, $\beta\ge1$ and the torus of side
$2n+4$, $\BP(n,\varpi)$ implies $\varpi\le C_Nn^4(1+\log\beta)/\beta$
\cite{BBLRS26}: on a small torus at weak coupling the Polyakov loops order, and
the cube dynamics can change their centre charge only through the curvature,
which is small. Along the tuning $n_k=\ell/a_k$ grows like a power of $1/a_k$
while $\beta_k$ grows like $\log(1/a_k)$, so this obstruction is consistent with
(D) but excludes every argument at fixed cube size.

\emph{(c) The group must be non-abelian.} For $G=\U(1)$ with the Villain
action at large $\beta$ the theory has a massless photon \cite{FS82}, and the
cube condition fails at every cube size, since by Theorem~\ref{thm:decay} it
would make the field-strength correlations summable \cite{BBLRS26}. Every proof
of (D) must therefore use the non-abelian structure of $G$, exactly as a proof
of the mass gap must.

\emph{The two ends.} Ba{\l}aban's renormalisation group
\cite{Bal84a,Bal85a,Bal89b,Dim13} controls the effective actions of the lattice
theory from the scale $a$ up to scales at which the running coupling is still
small, uniformly in the lattice spacing, in finite volume. The theorems of
Secs.~\ref{sec:curvature}--\ref{sec:bpstrong} control the theory when the
coupling is large: they need $N\beta$ small, that is, $g$ large. The running
coupling grows towards the infrared, so a proof of (D) would follow the
renormalisation group from the weak-coupling regime, where Ba{\l}aban's
estimates hold, through the regime of coupling of order one, to a scale
$\ell$ at which the effective theory satisfies a strong-coupling criterion of
the kind proved here, and would then transport the resulting cube estimate
back to the fine lattice. Neither the middle step nor the final transport is
available: the effective actions are not of Wilson form, Ba{\l}aban's analysis
stops before the coupling is of order one, and no known criterion accepts the
effective theory at that scale. This gap between the two ends is the Yang--Mills
mass gap problem in the form isolated here.

Statement (E) is the construction problem. Ba{\l}aban's ultraviolet stability
bounds give compactness of suitably smeared observables in finite volume, but
the infinite-volume limit, Euclidean invariance and nontriviality of the limit
are open \cite{JW06}. Given (E), Theorem~\ref{thm:B} reduces the mass gap to
(D) with the explicit constant $\Delta\ge16\varpi/(625e^2\ell)$.

\section{Conclusion}\label{sec:conclusion}

Lattice Yang--Mills theory in four dimensions has a mass gap with the explicit
lower bound $m_G(\beta)$ for every compact simple gauge group in the range
$N\beta<\frac2{29}\kappa_G$, and for $\SU(2)$ up to $\beta_W=\frac13$. A single
spectral gap $\varpi$ of the dynamics that resamples cubes of a fixed physical
size $\ell$, uniform along the continuum tuning, gives every nontrivial
continuum limit the mass gap $16\varpi/(625e^2\ell)$. That spectral gap holds
at strong coupling, for $\SU(2)$ up to $\beta_W=\frac5{18}$ and for $\SU(3)$ up
to $\beta_W=0.45$, and for massive free fields; it fails for massless fields,
at fixed cube size and in the abelian Coulomb phase, and along the weak-coupling
tuning required by the continuum problem it is not proved. The Clay problem
remains open; what this paper contributes is an explicit strong-coupling
theory for every gauge group and a reduction of the continuum gap to one
estimate on one scale.

\section*{Acknowledgements}

We thank I.~J.~Sessford for sharing his Lean development~\cite{Ses26}. The
von Mises--Fisher modulus and the Hessian bound without gauge fixing are those
of~\cite{BB26a}; the line field, the curvature bound in the gauge, the cube
condition and its consequences follow the longer account~\cite{BBLR26,BBLRS26}.
The constructive framework is that of Osterwalder, Schrader and Seiler, the
curvature criterion that of Bakry and \'Emery, and the uniqueness criterion
that of Dobrushin.

\section*{Declarations}

\emph{Data availability.} No data were generated. The longer account of this
work, with further results and finite checks, is archived at
\href{https://doi.org/10.5281/zenodo.24085959}{doi:10.5281/zenodo.24085959}.

\emph{Funding.} This research received no external funding.

\emph{Conflict of interest.} The authors declare no conflict of interest.

\emph{Ethics.} Not applicable: the work involves no human participants,
animals or personal data.

\appendix

\section{The line field}\label{app:table}

Table~\ref{tab:sigma} lists, for each residue class of $x$ modulo two, the
direction $\sigma(x)$ and the values $\pi_x(\mu)=\sigma(x+e_\mu)$. In every row
$\pi_x$ fixes $\sigma(x)$ and is a three-cycle on the other directions.

\begin{table}[h]
\caption{The line field: $x$ modulo two (digits $x_0x_1x_2x_3$), $\sigma(x)$, and
$\sigma(x+e_\mu)$ for $\mu=0,1,2,3$.}
\label{tab:sigma}
\begin{tabular}{cccccc|cccccc}
\toprule
$x$&$\sigma$&$0$&$1$&$2$&$3$&$x$&$\sigma$&$0$&$1$&$2$&$3$\\
\midrule
0000 & 0 & 0 & 2 & 3 & 1 & 0001 & 1 & 2 & 1 & 3 & 0 \\
1000 & 0 & 0 & 3 & 1 & 2 & 1001 & 2 & 1 & 3 & 2 & 0 \\
0100 & 2 & 3 & 0 & 2 & 1 & 0101 & 1 & 3 & 1 & 0 & 2 \\
1100 & 3 & 2 & 0 & 1 & 3 & 1101 & 3 & 1 & 2 & 0 & 3 \\
0010 & 3 & 1 & 2 & 0 & 3 & 0011 & 3 & 2 & 0 & 1 & 3 \\
1010 & 1 & 3 & 1 & 0 & 2 & 1011 & 2 & 3 & 0 & 2 & 1 \\
0110 & 2 & 1 & 3 & 2 & 0 & 0111 & 0 & 0 & 3 & 1 & 2 \\
1110 & 1 & 2 & 1 & 3 & 0 & 1111 & 0 & 0 & 2 & 3 & 1 \\
\bottomrule
\end{tabular}
\end{table}

\section{The blocks of Lemma~\ref{lem:local}}\label{app:blocks}

\emph{Diagonal blocks.} Let $w_a=e^{i\phi}$. On the diagonal part
$f_p=-\cos\phi\,(y_1+y_2-y_3-y_4)^2$ and $c_p=2(y_1y_4-y_1y_2+y_2y_3-y_3y_4)$, so
\[
\mathcal F=\sum_i\alpha_iy_i^2+\cos\phi\,(y_1+y_2-y_3-y_4)^2-\tfrac95(y_1y_4-y_1y_2
+y_2y_3-y_3y_4),
\]
with $y_k=0$, $\alpha_{k\pm1}=\frac{21}{10}$, $\alpha_{k+2}=\frac54$. It is affine in
$\cos\phi$, so positivity at $\cos\phi=\pm1$ suffices. For $k=1$ and $\cos\phi=1$
the matrix in $(y_2,y_3,y_4)$ is
\[
\begin{pmatrix}\frac{31}{10}&-\frac{19}{10}&-1\\[2pt]-\frac{19}{10}&\frac94&
\frac{19}{10}\\[2pt]-1&\frac{19}{10}&\frac{31}{10}\end{pmatrix},
\]
with leading principal minors $\frac{31}{10}$, $\frac{673}{200}$,
$\frac{8421}{2000}$. In general the leading principal minors are
$\frac{31}{10},\frac{673}{200},\frac{8421}{2000}$ (fixed link $e_1$ or $e_4$),
$\frac{31}{10},\frac{861}{100},\frac{8421}{2000}$ ($e_2$) and
$\frac94,\frac{673}{200},\frac{8421}{2000}$ ($e_3$) at $\cos\phi=1$, and
$\frac{11}{10},\frac{53}{200},\frac{21}{2000}$ ($e_1$, $e_4$),
$\frac{11}{10},\frac{21}{100},\frac{21}{2000}$ ($e_2$) and
$\frac14,\frac{53}{200},\frac{21}{2000}$ ($e_3$) at $\cos\phi=-1$; all are
positive.

\emph{Pair blocks.} Let $w_a=e^{i(\chi+\vartheta)}$ and $w_b=e^{i(\chi-\vartheta)}$. The
terms of $f_p$ carry $w_a+\bar w_b=2\cos\chi\,e^{i\vartheta}$ or its conjugate, and
$c_p$ carries $w_a\bar w_b=e^{2i\vartheta}$, so the Hermitian matrix $M(m,\vartheta)$ of
$\mathcal F$, $m=\cos\chi$, is affine in $m$ with $M(-1,\vartheta)=M(1,\vartheta+\pi)$,
and $m=1$ suffices. For the fixed link $e_2$, in the variables $(z_1,z_3,z_4)$
and with $c=\cos\vartheta$,
\[
M=\begin{pmatrix}\frac{21}5+2c&-2e^{i\vartheta}&-\frac95-2e^{i\vartheta}\\[2pt]
-2e^{-i\vartheta}&\frac{21}5+2c&\frac95+2e^{-i\vartheta}\\[2pt]
-\frac95-2e^{-i\vartheta}&\frac95+2e^{i\vartheta}&\frac52+2c\end{pmatrix}.
\]
With $u=\cos^2(\vartheta/2)\in[0,1]$ its leading principal minors are
\[
\tfrac{20u+11}5,\qquad\tfrac{(20u+1)(20u+21)}{25},\qquad
\tfrac{400u(40u^2-23u+4)+21}{250},
\]
and $40u^2-23u+4$ has negative discriminant $529-640$. The other positions give
the same determinant and the leading minors $\frac{20u+11}5$ or $\frac{8u+1}2$
and $\frac{800u^2-180u+53}{50}$, whose discriminant $32400-169600$ is negative.
Every block is positive definite.
